\documentclass[10pt,a4paper]{amsart}
\usepackage[margin=19mm]{geometry}
\usepackage{amsmath,amssymb,mathtools}
\usepackage[scr=rsfs,cal=euler]{mathalfa}
\usepackage{xfrac}
\usepackage[unicode,hidelinks,bookmarks=false]{hyperref}
\newcommand{\ie}{i.e.\ }
\newcommand{\cf}{cf.\ }
\newcommand{\resp}{respectively\ }
\newcommand{\Subsection}{\subsection}
\providecommand{\nobreakdash}{\nobreak\hbox{-}}
\setlength{\emergencystretch}{2em}
\multlinegap0pt
\usepackage{amssymb}
\newtheorem{lemma}{Lemma}
\DeclareMathOperator{\Ind}{Ind}
\def\ind{\mathop{\mathpalette\Ind{}}}
\def\nind{\mathop{\mathpalette\notind{}}}
\def\notind#1#2{#1\setbox0=\hbox{$#1x$}\kern\wd0
\hbox to 0pt{\mathchardef\nn=12854\hss$#1\nn$\kern1.4\wd0\hss}
\hbox to 0pt{\hss$#1\mid$\hss}\lower.9\ht0 \hbox to 0pt{\hss$#1\smile$\hss}\kern\wd0}
\DeclareMathOperator{\tp}{tp}
\title{Lascar Rank in Expansions}
\author{Michael Lange}
\date{Source: arXiv:2609.14698v1 (CC BY 4.0)}
\begin{document}
\maketitle

\noindent\emph{Source and licence.} Lascar Rank in Expansions. Version arXiv:2609.14698v1 (CC BY 4.0), distributed by arXiv under the Creative Commons Attribution 4.0 International licence, http://creativecommons.org/licenses/by/4.0/, as recorded in the arXiv licence metadata for this exact version and shown on the abstract page https://arxiv.org/abs/2609.14698v1. Full licence notice: \texttt{licenses/arxiv-2609.14698-CC-BY-4.0.txt}.\par\smallskip

\noindent\textit{Editorial scope.} Counted unit: the article's lemma EvenBetterBuechler (source0.tex lines 184--186). The article's complete original proof of it is delivered with it (source0.tex lines 188--218, one \texttt{proof} environment of 31 lines). No mathematics is written, repaired or re-derived: the statement and the proof are the article's own lines, in the article's own notation, and no retained line is substituted or re-flowed.  No further source-local context is needed: every cross-reference of every retained line resolves inside the two delivered spans.  Nothing was omitted from any retained span, so the package carries no omission record.  No retained line contains a citation, so the wrapper carries no bibliography and no bibliography entry.  No retained line contains a figure environment, an image-inclusion command or drawing code, so no image is delivered and the package declares no asset dependency.  Editorial formatting only, in addition: an AMS \texttt{amsart} preamble that restates the article's own declarations and macros used by the retained lines, namely the environments \texttt{lemma} on separate counters, together with the standard packages those lines need.  The retained environments print in this document as Lemma 1 in order of appearance, whereas the article prints the same items with the numbering of its own full document.  The complete native source of the article as distributed in the arXiv e-print for this exact version is bundled unchanged as \texttt{sources/210349/source0.tex}.
\begin{lemma}\label{EvenBetterBuechler}
    Let $p\in S(A)$ be a complete type and $q\in S(N)$ an extension of $p$, with $N$ being $|A|^+$-saturated. Suppose also that $\mathrm{SU}(q)+\omega^\alpha\leq\mathrm{SU}(p)$. Then there is an intermediate extension $p\subseteq r\subseteq q$ with $\mathrm{SU}(r)\approx_\alpha \mathrm{SU}(q)+\omega^\alpha$. %In particular $r\in S(Ab)$ for some finite tuple $b\in N$.
\end{lemma}

\begin{proof}
    By elimination of hyperimaginaries, we have $\mathrm{Cb}(q)\subseteq N$ and by supersimplicity, some finite tuple $c\in \mathrm{Cb}(q)$ such that $q$ does not fork over $c$. If $a$ is any realization of $q$, we have $\mathrm{SU}(a/Ac)+\omega^\alpha\leq\mathrm{SU}(a/A)$. Then by quantitative forking symmetry we have $\mathrm{SU}(c/A)\geq \mathrm{SU}(c/Aa)+\omega^\alpha$. In particular, we have that $\mathrm{SU}(c/A)\geq \omega^\alpha$ and so there exists a finite tuple $b$ such that $\mathrm{SU}(c/Ab)=\omega^\alpha$. By saturation, it may be assumed that $b\in N$. Let $a\models q$ so in particular $a\ind_{Ac}b$.

We claim that $a\nind_{Ab} c$. Indeed, suppose not: then by transitivity we have that $a\ind_{Ab} N$ and so $\mathrm{Cb}(a/N)\subseteq \mathrm{acl}(Ab)$; in particular, $c\in \mathrm{acl}(Ab)$ which contradicts $\mathrm{SU}(c/Ab)=\omega^\alpha$.

Now by symmetry, $c\nind_{Ab} a$ which in particular means $\mathrm{SU}(c/Aab)<\omega^\alpha$.

Let us consider two ways of bounding $\mathrm{SU}(ac/Ab)$ using the Lascar inequalities. First:

\[\mathrm{SU}(c/Aab)+\mathrm{SU}(a/Ab)\leq\mathrm{SU}(ac/Ab)\leq\mathrm{SU}(c/Aab)\oplus\mathrm{SU}(a/Ab)\]

This inequality and the fact that $\mathrm{SU}(c/Aab)<\omega^\alpha$ yield:

\[\mathrm{SU}(a/Ab)\approx_\alpha\mathrm{SU}(ac/Ab)\]

Now the other:

\[\mathrm{SU}(a/Abc)+\mathrm{SU}(c/Ab)\leq\mathrm{SU}(ac/Ab)\leq\mathrm{SU}(a/Abc)\oplus\mathrm{SU}(c/Ab)\]

The assumption that $\mathrm{SU}(c/Ab)=\omega^\alpha$ implies that 
\[\mathrm{SU}(a/Abc)\oplus\mathrm{SU}(c/Ab)\approx_\alpha\mathrm{SU}(a/Abc)+\mathrm{SU}(c/Ab)=\mathrm{SU}(q)+\omega^\alpha\]

and so the previous inequality yields:
\[\mathrm{SU}(ac/Ab)\approx_\alpha\mathrm{SU}(q)+\omega^\alpha\]

which together with our first $\approx_\alpha$ equivalence gives:

\[\mathrm{SU}(a/Ab)\approx_\alpha\mathrm{SU}(q)+\omega^\alpha\]

Hence we may take $r=\tp(a/Ab)$ to conclude the claim.
\end{proof}

\end{document}
